
Ohta (2026) shows that price clustering on the Tokyo Stock Exchange arises
because noise traders leave limit orders standing at round prices and faster
participants take them with large market orders, and closes by proposing that
clustering measures could therefore serve as an observable proxy for
noise-trader activity --- leaving the relationship between that activity, price
formation and liquidity as future work. This study carries out that analysis on
the 2025 tape, three years beyond the paper's sample, using 77,592 stock-days
across 1,046 stocks and 78 trading days.

The measures replicate. Round prices carry 14.95\% of large-trade volume
against 10.77\% of small-trade volume, close to the paper's 14.8\% and
11.1\%, and the large-minus-small ordering its mechanism predicts holds with
high confidence. Trades at round prices move the midquote by about
0.43 basis points more than other trades on the buy side, against
the paper's 1.26.

Two measures the paper's data could not support are added. The share of visible
ten-level resting depth sitting at round prices averages 13.21\% against a
uniform benchmark of 10\%, observing the stale inventory directly rather than
after it has been executed against; and limit-order submission and cancellation
shares are inferred from ladder deltas, reproducing the paper's published
magnitudes without having been calibrated to them.

On the open question --- whether noise-trader activity marks a more liquid market
or a more expensive one --- the panel gives an answer, reported with the caution
a four-month sample deserves. A placebo that rebuilds the measure on each of the other
nine last digits separates the round-price result from a mechanical artefact of
how the measure is constructed. A strategy demonstration is included to exercise
the machinery end to end; its informative output is the cost accounting, not the
return.
